\section{Projection tools}\label{r:sec:projections}

Throughout this section, $D$ is a unital $C^*$-algebra of real rank zero
with $\T(D)\ne\varnothing$.  The trace estimates below hold simultaneously
for all $\tau\in\T(D)$.  The finite-set construction needs subprojections
with prescribed values on finitely many trace tests, followed by a way to
move them off previously occupied projections.  Scalar compressions and
proportional splitting provide the first step; the last lemma controls
the trace loss and displacement in the second.

\begin{lemma}[Scalar compressions]\label{r:proj:compress}
Every projection corner $rDr$ has real rank zero.  Moreover, given a
projection $r\in D$, a finite set $\mathcal F\subset D$, and $\gamma>0$,
there is a finite family of nonzero orthogonal projections $q_1,\ldots,q_m$
with sum $r$ such that
\begin{equation}\label{r:proj:compress-estimate}
  \dist(q_j hq_j,\C q_j)<\gamma
  \qquad(h\in\mathcal F,\ 1\le j\le m).
\end{equation}
The family is empty when $r=0$.
\end{lemma}

\begin{proof}
Let $r\ne0$ and $h=h^*\in rDr$.  Choose $R>\norm{h}+2$ and finite-spectrum
self-adjoints $a_n\in D$ converging in norm to $h+Rr$.  The spectrum of the
limit is contained in $\{0\}\cup(2,\infty)$, so the spectral projections
$f_n=1_{(1,\infty)}(a_n)$ converge in norm to $r$.  For all sufficiently
large $n$, the element
\[
  r f_n+(1-r)(1-f_n)
\]
is invertible.  Its unitary polar part $u_n$ satisfies
$u_n f_n u_n^*=r$ and $u_n\to1$ in norm.  Since $a_n$ commutes with $f_n$,
the compression $r u_n a_n u_n^*r$ has finite spectrum in $rDr$.
Subtracting $Rr$ gives finite-spectrum approximations to $h$ in that
corner.

For the second assertion, first replace $\mathcal F$ by its real and
imaginary parts.  Compress the first self-adjoint test to $rDr$, approximate
it by a finite-spectrum self-adjoint, and partition $r$ by the latter's
spectral projections.  Repeat for each subsequent test in every current
corner.  An estimate
$\norm{q hq-\lambda q}<\delta$ persists on every subprojection of $q$.
Thus this finite refinement procedure preserves all earlier estimates.
Using tolerance less than $\gamma/2$ for the real and imaginary parts
proves \eqref{r:proj:compress-estimate}; zero pieces are discarded.
\end{proof}

The next fact concerns a simple unital algebra alone. It uses neither
real rank zero nor nuclearity nor comparison, and will also supply the
small hereditary support for the rank approximation in
Section~\ref{u:sec:ranks}.

\begin{lemma}[Full positive elements with small tracial support]
\label{u:lem:thin-full}
Let $A$ be simple, unital and infinite dimensional, with
$T(A)\ne\varnothing$.  For every $\delta>0$ there is a full positive element
$e\in A$ with $\norm e=1$ such that
\[
 \tau(\supp(e))<\delta\qquad(\tau\in T(A)).
\]
Here supports and traces are taken in each tracial von Neumann representation.
\end{lemma}

\begin{proof}
First we obtain arbitrarily many mutually orthogonal nonzero positive
elements.  Choose a maximal abelian unital $C^*$-subalgebra $E\subset A$.
Suppose that $E$ were finite dimensional, with minimal projections
$p_1,\ldots,p_m$ summing to $1$.  Every selfadjoint element of $p_iAp_i$
commutes with $E$, so maximality implies $p_iAp_i=\C p_i$.  If
$0\ne x\in p_iAp_j$, then
\[
 x^*x=\norm{x}^2p_j,\qquad xx^*=\norm{x}^2p_i.
\]
Thus $u=x/\norm{x}$ satisfies $u^*u=p_j$ and $uu^*=p_i$.  For any
$y\in p_iAp_j$, the element $u^*y$ is scalar on $p_j$, and
$y=u(u^*y)$.  Consequently each $p_iAp_j$ has dimension at most one.
The decomposition $A=\sum_{i,j}p_iAp_j$ would make $A$ finite
dimensional, a contradiction.  Hence $E\cong C(X)$ for an infinite
compact Hausdorff space $X$.  Given $k\in\N$, choose $k$ distinct points
of $X$, pairwise disjoint neighborhoods, and nonzero positive continuous
functions supported in those neighborhoods.  Their images give mutually
orthogonal nonzero positive elements $c_1,\ldots,c_k\in A$.

For two nonzero positive elements $c,d\in A$, simplicity gives $cAd\ne0$.
Indeed, if $cAd=0$, then $c$ annihilates the ideal generated by $d$;
this ideal is $A$, which would force $c=0$.  Choose $y\in A$ such that
$z=cyd\ne0$.  Then $z^*z\in\overline{dAd}$ and $zz^*\in\overline{cAc}$.
In each tracial von Neumann representation, the polar decomposition of
$z$ gives a partial isometry with
initial projection $\supp(z^*z)$ and final projection $\supp(zz^*)$.
The normal tracial state therefore gives
\[
 \tau(\supp(z^*z))=\tau(\supp(zz^*))
 \le \min\{\tau(\supp(c)),\tau(\supp(d))\}
 \qquad(\tau\in T(A)).
\]
The elements $y,z$ are chosen independently of $\tau$; the support
comparison holds in every tracial representation.
Start with $c=c_1$. For $j=2,\ldots,k$, apply the preceding construction
with $d=c_j$ and replace $c$ by $zz^*$. The support trace never increases
and is at most the support trace of $c_j$. The final nonzero positive
element $e=c$ therefore has support trace at most that of each $c_i$,
at every trace.  Since the support
projections of the $c_i$ are orthogonal,
\[
 k\,\tau(\supp(e))\le\sum_{i=1}^k\tau(\supp(c_i))\le1
 \qquad(\tau\in T(A)).
\]
Choose $k$ with $1/k<\delta$ and rescale $e$ to have norm one, which
does not change its support.  It is full because $A$ is simple.
\end{proof}

We return to the ambient algebra $D$. Real rank zero will produce a
projection by a spectral cut; the norm-one normalization will let us prove
that this projection is full by testing proper quotients of $D$.

\begin{lemma}[Small full projections]\label{r:proj:small}
Suppose that $D$ contains a unital simple infinite-dimensional
$C^*$-subalgebra $A$.  For every $\delta>0$, there is a full projection
$g\in D$ such that
\[
  \sup_{\tau\in\T(D)}\tau(g)<\delta.
\]
\end{lemma}

\begin{proof}
The restriction to $A$ of any trace in $\T(D)$ is a tracial state,
so $\T(A)\ne\varnothing$. Given a positive integer $n$, apply
Lemma~\ref{u:lem:thin-full} to obtain $b\in A_+$ with $\norm b=1$ and
\begin{equation}\label{r:proj:small-positive}
 \rho(b)\le\rho(\supp(b))<\frac1n
 \qquad(\rho\in\T(A)).
\end{equation}

Choose a finite-spectrum self-adjoint $a\in D$ with
$\norm{a-b}<1/8$, and set $g=1_{[1/2,\infty)}(a)$.  Then
$gbg\ge(3/8)g$, so restriction of any $\tau\in\T(D)$ to $A$ gives
\begin{equation}\label{r:proj:small-projection-bound}
  \tau(g)\le\frac83\tau(b)\le\frac{8}{3n}.
\end{equation}
We claim that $g$ is full in $D$.  Otherwise, let $\pi:D\to D/I$ be the quotient
by the proper closed ideal generated by $g$.  The inequalities
$a\ge-1/8$ and $a(1-g)\le(1/2)(1-g)$ imply
$\norm{\pi(a)}\le1/2$, and hence $\norm{\pi(b)}\le5/8$.
On the other hand, $\pi|_A$ is a unital homomorphism from the simple
algebra $A$ to a nonzero unital algebra.  It is injective and therefore
isometric, contradicting $\norm{b}=1$.  Thus $g$ is full, and choosing
$n$ with $8/(3n)<\delta$ finishes the proof.
\end{proof}

We will use an elementary splitting procedure with different target
projections.  Zero projections are allowed among the resulting pieces.

\begin{lemma}[Splitting a column]\label{r:proj:columns}
Let $f,g_1,\ldots,g_m$ be projections in $D$, where $m\ge1$.
Suppose that $v_1,\ldots,v_m\in D$ satisfy
\[
  v_i=g_i v_i f,
  \qquad \sum_{i=1}^m v_i^*v_i=f.
\]
There are orthogonal projections $e_1,\ldots,e_m$ with sum $f$ and
$e_i\precsim g_i$ for every $i$.
\end{lemma}

\begin{proof}
If $f=0$, take all $e_i=0$.  We argue by induction on $m$; for $m=1$,
$v_1$ itself implements $f\precsim g_1$.  Suppose that $m>1$.
Apply Lemma~\ref{r:proj:compress} to choose a finite-spectrum self-adjoint
$a\in fDf$ within $1/8$ of $v_1^*v_1$.  Set
\[
  e_1=1_{[1/2,\infty)}(a),\qquad r=f-e_1.
\]
The inequality $e_1v_1^*v_1e_1\ge(3/8)e_1$ shows that, when $e_1\ne0$,
\[
  v_1e_1(e_1v_1^*v_1e_1)^{-1/2}
\]
is a partial isometry with initial projection $e_1$ and final projection
below $g_1$.  On the remaining corner,
\[
  Q:=\sum_{i=2}^m r v_i^*v_i r
    =r-rv_1^*v_1r\ge\frac38r.
\]
If $r=0$, take $e_2=\cdots=e_m=0$.  Otherwise define
$w_i=v_i r Q^{-1/2}$ for $i=2,\ldots,m$, with the inverse taken in $rDr$.
These elements satisfy $w_i=g_iw_i r$ and
$\sum_{i=2}^m w_i^*w_i=r$.  The induction hypothesis partitions $r$ as
required.
\end{proof}

Approximate projection halving in non-elementary simple real-rank-zero
algebras was developed by Zhang \cite[Theorem~1.1(i)]{Zhang1991}.
More directly, Zhang's ordered-decomposition lemma
\cite[Lemma~1.1]{Zhang1990}, as used in Vaccaro's divisibility
argument~\cite[Lemma~2.1, version~1]{Vac26}, applies without
simplicity: a projection can be
partitioned into finitely many pieces ordered by subequivalence and each
subequivalent to any given full projection. We give this projection
construction below and then use a common rounding rule to obtain
proportional trace splitting uniformly over all traces.

\begin{proposition}[Uniform proportional subprojections]\label{r:proj:split}
Assume that for every $\delta>0$ there is a full projection $g\in D$ with
$\sup_{\tau\in\T(D)}\tau(g)<\delta$.
For every projection $f\in D$, every $s\in[0,1]$, and every $\eps>0$,
there is a projection $e\le f$ such that
\begin{equation}\label{r:proj:split-estimate}
  \sup_{\tau\in\T(D)}|\tau(e)-s\tau(f)|<\eps.
\end{equation}
\end{proposition}

\begin{proof}
The case $f=0$ is immediate.  Choose a full projection $g$ with
$\sup_\tau\tau(g)<\eps$.  Fullness provides a finite sum
$\sum_i x_i g y_i$ within distance less than one of $1$.
Write $X=[x_1\ \cdots\ x_m]$ and let $Y$ be the column with entries
$gy_i$.  Then $XY$ is invertible, so $Y$ has a bounded left inverse
$L=(XY)^{-1}X$.  In particular,
\[
  1=Y^*L^*LY\le\norm{L}^2Y^*Y,
\]
and $Y^*Y$ is invertible.  The entries of
$Y(Y^*Y)^{-1/2}f$ form a column as in Lemma~\ref{r:proj:columns}, with every
target projection equal to $g$.  We obtain an orthogonal family with
sum $f$, each member subequivalent to $g$.

We next replace this family by one, say $s_1,\ldots,s_m$, satisfying
\begin{equation}\label{r:proj:ordered-partition}
  s_1\precsim s_2\precsim\cdots\precsim s_m\precsim g.
\end{equation}
This ordering makes the trace weights increase in the same order for
every trace, so that one choice of rounding coefficients will work
simultaneously for all of them.
Here is the pair operation that makes this possible.  Suppose that $x,y$
are orthogonal pieces, both subequivalent to a common projection $h$.
Choose partial isometries $w,u$ with
\[
  w^*w=x,\quad ww^*=x'\le h,
  \qquad u^*u=y,\quad uu^*\le h.
\]
Since $u^*x'u+u^*(h-x')u=y$, Lemma~\ref{r:proj:columns} applied to the
column with entries $x'u,(h-x')u$ gives $y=k+l$ and partial isometries
$\alpha,\beta$ with
\[
  \alpha^*\alpha=k,\quad \alpha\alpha^*\le x',
  \qquad \beta^*\beta=l,\quad \beta\beta^*\le h-x'.
\]
Set $\xi=w^*\alpha+l$ and $\zeta=w+\beta$.  Then
\begin{align*}
  \xi^*\xi&=y,&
  \xi\xi^*&=w^*\alpha\alpha^*w+l\le x+l,\\
  \zeta^*\zeta&=x+l,&
  \zeta\zeta^*&=x'+\beta\beta^*\le h.
\end{align*}
The cross terms vanish by orthogonality of the initial and final
components: $x\perp l$ and $x'\perp(h-x')$.
Thus replacing $x,y$ by $k,x+l$
preserves their sum and orthogonality, and
\[
  k\precsim x\le x+l,
  \qquad x,y\precsim x+l\precsim h.
\]
After each adjacent comparison, the new right-hand piece dominates
every piece already passed: it contains the previous moving piece.
Thus a left-to-right pass with common bound $g$ leaves a last piece
dominating all earlier pieces.  Freeze that last
piece and repeat on the shorter list, now using it as the common bound.
Induction gives \eqref{r:proj:ordered-partition}.  All operations take place
inside $D$ and preserve the orthogonal sum $f$.

Set
\[
  c_i=\lfloor si\rfloor-\lfloor s(i-1)\rfloor\in\{0,1\},
  \qquad e=\sum_{i=1}^m c_i s_i.
\]
For a fixed $\tau\in\T(D)$, put $t_i=\tau(s_i)$ and
$D_j=\lfloor sj\rfloor-sj$.  Then $t_1\le\cdots\le t_m\le\tau(g)$ and
$-1\le D_j\le0$.  Summation by parts gives
\[
  \tau(e)-s\tau(f)
   =D_m t_m+\sum_{i=1}^{m-1}D_i(t_i-t_{i+1}).
\]
The first term lies in $[-t_m,0]$, while the sum lies in
$[0,t_m-t_1]$.  It follows that
$|\tau(e)-s\tau(f)|\le t_m\le\tau(g)$, proving
\eqref{r:proj:split-estimate} uniformly in $\tau$.
\end{proof}

\begin{corollary}[Splitting finite linear tests]\label{r:proj:tests}
Under the small-full-projection hypothesis of Proposition~\ref{r:proj:split},
let $q\in D$ be a projection, $s\in[0,1]$, and $\mathcal F\subset D$ a
finite set.  For every $\eps>0$, there is a projection $e\le q$ such that
\begin{equation}\label{r:proj:tests-estimate}
  |\tau(eh)-s\tau(qh)|<\eps
  \qquad(h\in\mathcal F,\ \tau\in\T(D)).
\end{equation}
\end{corollary}

\begin{proof}
The assertion is immediate if $q=0$ or $\mathcal F$ is empty.  Choose
$\gamma>0$ with $2\gamma<\eps/2$ and use
Lemma~\ref{r:proj:compress} to write $q=\sum_{j=1}^m q_j$ with
\[
  \norm{q_j hq_j-\lambda_{j,h}q_j}<\gamma
  \qquad(h\in\mathcal F)
\]
for suitable scalars $\lambda_{j,h}$.  These scalars satisfy
$|\lambda_{j,h}|\le\norm{h}+\gamma$.
Put $M=\max_{h\in\mathcal F}\norm{h}+\gamma$ and choose
$\delta>0$ with $mM\delta<\eps/2$.
By Proposition~\ref{r:proj:split}, choose $e_j\le q_j$ so that
$|\tau(e_j)-s\tau(q_j)|<\delta$ for every $\tau$, and set
$e=\sum_j e_j$.  Traciality removes all off-diagonal compressions and gives
\begin{align*}
  |\tau(eh)-s\tau(qh)|
  &\le \gamma\bigl(\tau(e)+s\tau(q)\bigr)
       +\sum_{j=1}^m |\lambda_{j,h}|
          |\tau(e_j)-s\tau(q_j)| \\
  &\le 2\gamma+mM\delta<\eps.
\end{align*}
\end{proof}

The last tool moves a projection off an occupied corner.  Its linear
trace estimate retains the scale of the projection being moved.

\begin{lemma}[Orthogonal insertion]\label{r:proj:move}
For projections $z,y\in D$, there is a projection $t\le1-z$ such that
\begin{align}
  \tau(y)-3\tau(yz)&\le\tau(t)\le\tau(y),
       \label{r:proj:move-mass}\\
  \normtwo{t-y}{\tau}^{2}&\le5\tau(yz)
       \label{r:proj:move-distance}
\end{align}
for every $\tau\in\T(D)$.  For every contraction $x\in D$, it also
satisfies
\begin{equation}\label{r:proj:move-linear}
  |\tau((t-y)x)|
   \le2\sqrt{\tau(y)}\,\normtwo{t-y}{\tau}.
\end{equation}
\end{lemma}

\begin{proof}
If $y=0$, take $t=0$.  Otherwise choose, in $yDy$, a finite-spectrum
self-adjoint $a$ with $\norm{a-yzy}<1/8$, and let
$f=1_{(-\infty,1/2)}(a)$.  Writing $d=y-f$, we have
\[
  d(yzy)d\ge\frac38d,
  \qquad f(yzy)f\le\frac58f.
\]
Since $yzy\ge0$, this implies
\begin{equation}\label{r:proj:discarded-mass}
  \tau(y-f)\le\frac83\tau(yz),
  \qquad A_f:=f(1-z)f\ge\frac38f.
\end{equation}
If $f=0$, put $t=0$.  Otherwise the element
\[
  v=(1-z)f A_f^{-1/2}
\]
is a partial isometry with $v^*v=f$; set $t=vv^*\le1-z$.
In both cases,
\[
  \tau(t)=\tau(f),\qquad ftf=f(1-z)f.
\]
The mass bounds \eqref{r:proj:move-mass} follow from
\eqref{r:proj:discarded-mass}.  As $f\le y$,
\[
  \tau(ty)\ge\tau(tf)=\tau(f)-\tau(fz).
\]
Consequently,
\begin{align*}
  \normtwo{t-y}{\tau}^{2}
   &=\tau(t)+\tau(y)-2\tau(ty)\\
   &\le\tau(y-f)+2\tau(fz)
    \le\left(\frac83+2\right)\tau(yz)
    \le5\tau(yz),
\end{align*}
which proves \eqref{r:proj:move-distance}.

Finally, decompose $t-y=t(t-y)+(1-t)(t-y)$.  The first summand has left
support below $t$, and the second equals $-(1-t)y$ and has right support
below $y$.  Tracial Cauchy--Schwarz therefore gives
\[
  \normone{t-y}{\tau}
   \le\bigl(\sqrt{\tau(t)}+\sqrt{\tau(y)}\bigr)
       \normtwo{t-y}{\tau}
   \le2\sqrt{\tau(y)}\,\normtwo{t-y}{\tau}.
\]
Trace-norm duality proves \eqref{r:proj:move-linear}.
\end{proof}
